\documentclass[10pt,a4paper]{article}
\usepackage[utf8]{inputenc}
\usepackage[italian]{babel}
\usepackage[left=2cm,right=2cm,top=2cm,bottom=2cm]{geometry}
\usepackage{titlesec}
\usepackage{xcolor}
\usepackage{hyperref}
\usepackage{array}

% Definizione palette colori (ispirata a Linux Mint)
\definecolor{mintdark}{HTML}{2F5233}
\definecolor{mintlight}{HTML}{87A96B}
\definecolor{text}{HTML}{222222}

\titleformat{\section}{\large\bfseries\color{mintdark}}{}{0em}{}[\titlerule]
\titlespacing*{\section}{0pt}{14pt}{6pt}

\begin{document}
\color{text}
\pagestyle{empty}

% --- INTESTAZIONE CONTATTI (In alto) ---
\noindent
\textbf{E-mail:} \href{mailto:stefanoperis.github@gmail.com}{stefanoperis.github@gmail.com} \hfill 
\textbf{GitHub:} \href{https://github.com/TechnoKnight77}{https://github.com/TechnoKnight77}

\vspace{15pt}

% --- BLOCCO TITOLO, AUTORE E DATA ---
\begin{center}
    {\LARGE \bfseries \color{mintdark} Guida ai Comandi da Terminale di Linux (sistemi basati su Debian)} \\[0.4cm]
    {\large Stefano Peris} \\[0.2cm]
    {\small Ultima modifica: \today}
\end{center}

\vspace{15pt}

% --- NAVIGAZIONE ---
\section*{Navigazione nei File e nelle Cartelle}
\noindent
\begin{tabular}{@{}p{5cm}p{12cm}@{}}
\textbf{\texttt{pwd}} & Mostra la cartella corrente in cui ti trovi (Print Working Directory). \\
\textbf{\texttt{ls}} & Elenca i file e le cartelle presenti nella directory corrente. \\
\textbf{\texttt{ls -la}} & Elenca tutti i file, inclusi quelli nascosti, con dettagli su permessi e dimensioni. \\
\textbf{\texttt{cd [cartella]}} & Cambia la cartella corrente ed entra in quella specificata. \\
\textbf{\texttt{cd ..}} & Torna indietro di una cartella nella gerarchia del sistema. \\
\textbf{\texttt{cd \textasciitilde}} & Torna istantaneamente alla tua cartella Home personale.
\end{tabular}

% --- GESTIONE FILE ---
\section*{Gestione di File e Cartelle}
\noindent
\begin{tabular}{@{}p{5cm}p{12cm}@{}}
\textbf{\texttt{mkdir [nome\_cartella]}} & Crea una nuova cartella con il nome specificato. \\
\textbf{\texttt{touch [nome\_file]}} & Crea un file di testo vuoto o ne aggiorna la data di modifica. \\
\textbf{\texttt{cp [origine] [destinazione]}} & Copia un file o una cartella da un percorso a un altro. \\
\textbf{\texttt{mv [origine] [destinazione]}} & Sposta o rinomina un file o una cartella. \\
\textbf{\texttt{rm [nome\_file]}} & Elimina definitivamente il file specificato. \\
\textbf{\texttt{rm -rf [nome\_cartella]}} & Elimina forzatamente una cartella e tutto il suo contenuto (usare con cautela).
\end{tabular}

% --- GESTIONE PACCHETTI ---
\section*{Gestione dei Pacchetti e Aggiornamenti (APT)}
\noindent
\begin{tabular}{@{}p{5cm}p{12cm}@{}}
\textbf{\texttt{sudo apt update}} & Sincronizza l'indice dei pacchetti dai server per verificare gli aggiornamenti disponibili. \\
\textbf{\texttt{sudo apt upgrade}} & Installa gli aggiornamenti disponibili per tutti i software del sistema. \\
\textbf{\texttt{sudo apt install [pacchetto]}} & Scarica e installa un nuovo software dai repository ufficiali. \\
\textbf{\texttt{sudo apt remove [pacchetto]}} & Rimuove un software installato mantenendo i file di configurazione. \\
\textbf{\texttt{sudo apt autoremove}} & Rimuove automaticamente i pacchetti e le dipendenze non più necessari.
\end{tabular}

% --- INFORMAZIONI DI SISTEMA ---
\section*{Informazioni di Sistema e Monitoraggio}
\noindent
\begin{tabular}{@{}p{5cm}p{12cm}@{}}
\textbf{\texttt{df -h}} & Mostra lo spazio disponibile e utilizzato sui dischi in formato leggibile (GB, MB). \\
\textbf{\texttt{free -h}} & Mostra l'utilizzo della memoria RAM e della partizione di Swap del computer. \\
\textbf{\texttt{top}} & Mostra in tempo reale i processi attivi e le risorse hardware utilizzate. \\
\textbf{\texttt{uname -a}} & Stampa le informazioni dettagliate sul kernel Linux e sul sistema in uso.
\end{tabular}

% --- SCORCIATOIE ---
\section*{Scorciatoie Essenziali da Tastiera}
\noindent
\begin{tabular}{@{}p{5cm}p{12cm}@{}}
\textbf{\texttt{Tab}} & Autocompleta automaticamente i nomi dei comandi, dei file o delle cartelle. \\
\textbf{\texttt{Ctrl + C}} & Interrompe e blocca immediatamente il comando o il processo in esecuzione. \\
\textbf{\texttt{Ctrl + ALT + T}} & Scorciatoia di sistema su Linux Mint per aprire istantaneamente il Terminale. \\
\textbf{\texttt{clear}} & Pulisce interamente la schermata del terminale da tutto il testo precedente.
\end{tabular}

\end{document}
